%\documentclass[letterpaper,12pt]{article}
\documentclass[preview]{standalone}

\usepackage{tabularx}
\usepackage{amsmath}
\usepackage{graphicx}
\usepackage[margin=1in,letterpaper]{geometry}
\usepackage{cite}
\usepackage[final]{hyperref}
\usepackage{datatool}
\usepackage{caption}
\usepackage{float}
\usepackage{gensymb}

\hypersetup{
        colorlinks=true,
        linkcolor=blue,
        citecolor=blue,
        filecolor=magenta,
        urlcolor=blue
}
\begin{document}
 \setcounter{table}{4}
 \begin{table}[h!]
 \centering
 \label{tab:fits}
 \renewcommand{\arraystretch}{1.2}
 \begin{tabular}{|c|c|c|c|c|}
 \hline
 Noyau & $A$ | V & $\lambda$ & $\omega$ | rad/s & T | s \\
 \hline
 Acier     &
 $4.597 \pm 0.005$ &
 $263.6 \pm 0.4$ &
 $5362.9 \pm 0.4$ &
 $0.0011716 \pm 9 \cdot 10^{-8}$ \\

 Acier feuilleté     &
 $4.565 \pm 0.005$ &
 $187.4 \pm 0.3$ &
 $3381.8 \pm 0.3$ &
 $0.0018579 \pm 2 \cdot 10^{-7}$ \\

 Air     &
 $3.75 \pm 0.01$ &
 $1276 \pm 6$ &
 $7014 \pm 6$ &
 $0.0008958 \pm 8 \cdot 10^{-7}$ \\

 Aluminium    &
 $4.212 \pm 0.006$ &
 $1173 \pm 3$ &
 $8875 \pm 3$ &
 $0.0007079 \pm 2 \cdot 10^{-7}$ \\
 \hline
\end{tabular}
 \caption{Paramètres ajustés des oscillations avec la boite à décade et R = $10 \Omega$}
 \end{table}


\end{document}

